% ============================================================
% Round 07: Flow-GRPO Foundation
% ============================================================

\subsection{Round 07：Flow-GRPO：Flow 侧 RL 基础}

\subsubsection{Highlight}

\begin{conceptbox}
\textbf{本章相关脉络：}
\begin{itemize}[leftmargin=1.5em]
  \item \textbf{Flow-GRPO}：arXiv 2025-05-08，NeurIPS 2025，Training Flow Matching Models via Online RL。核心是把 deterministic ODE 转成可采样、可算 logprob 的 SDE，再用 GRPO 做 online RL。论文：\url{https://arxiv.org/abs/2505.05470}，代码：\url{https://github.com/yifan123/flow_grpo}
  \item \textbf{Flow-GRPO 官方实现}：2025-08-15 加入 Qwen-Image / Qwen-Image-Edit 支持；2025-08-04 加入 FLUX.1-Kontext-dev 图像编辑 pipeline。仓库：\url{https://github.com/yifan123/flow_grpo}
  \item \textbf{LeapAlign}：arXiv 2026-04-16，CVPR 2026，Post-Training Flow Matching Models at Any Generation Step by Building Two-Step Trajectories。它不是 policy gradient，而是把长 ODE 轨迹压缩成 two-step leap trajectory，让 reward gradient 可以稳定传到早期生成步骤。论文：\url{https://arxiv.org/abs/2604.15311}，项目页：\url{https://rockeycoss.github.io/leapalign/}
\end{itemize}
\end{conceptbox}

\subsubsection{本章要解决的问题}

上一章 DDPO 解决的是 diffusion model 的 RL：

\[
p_\theta(z_{t-1}\mid z_t,x,t)
\]

每一步本来就是随机转移，所以可以算：

\[
\log p_\theta(z_{t-1}\mid z_t,x,t)
\]

然后做 PPO / DDPO ratio。

Flow Matching 的问题更麻烦。它的采样默认是 ODE：

\[
\frac{dx_t}{dt}
=
v_\theta(x_t,t,x)
\]

离散采样时：

\[
x_{t-\Delta t}
=
x_t-\Delta t\,v_\theta(x_t,t,x)
\]

这是确定性的。

因此它默认没有一个像 diffusion 那样好用的：

\[
p_\theta(x_{t-\Delta t}\mid x_t,x,t)
\]

\begin{warnbox}
\textbf{Flow RL 的核心难点：} DDPO 可以直接把 diffusion reverse transition 当 policy；Flow Matching 的 ODE transition 默认只是一个点，不是一个普通概率分布，所以不能直接做 logprob ratio。
\end{warnbox}

\subsubsection{从 DDPO 到 Flow-GRPO 的概念迁移}

\begin{center}
{\small
\begin{tabular}{p{3cm}p{4.8cm}p{4.8cm}}
\hline
\textbf{概念} & \textbf{DDPO / Diffusion RL} & \textbf{Flow-GRPO} \\
\hline
生成过程
&
随机 denoising chain
&
默认确定性 ODE flow
\\
每步转移
&
$p_\theta(z_{t-1}\mid z_t,x,t)$
&
先把 ODE 转成 SDE，得到 $p_\theta(x_{t-\Delta t}\mid x_t,x,t)$
\\
Action
&
下一步 latent $z_{t-1}$
&
下一步 latent $x_{t-\Delta t}$
\\
Logprob
&
天然来自 Gaussian reverse transition
&
来自 ODE-to-SDE 后的 Gaussian transition
\\
Advantage
&
reward baseline / PPO advantage
&
group-relative advantage
\\
Ratio
&
新旧 denoising transition probability ratio
&
新旧 flow transition probability ratio
\\
\hline
\end{tabular}
}
\end{center}

\begin{summarybox}
\textbf{一句话：} Flow-GRPO 的目标不是重新发明 PPO，而是先让 flow transition 变成可采样、可算 logprob 的随机策略，然后把 GRPO 的 group-relative ratio loss 接上去。
\end{summarybox}

\subsubsection{第一步：把 ODE 写成一步均值}

按照本笔记前面的时间约定，$t=1$ 是噪声端，$t=0$ 是干净端。采样时从大 $t$ 往小 $t$ 走。

ODE 的 Euler 更新是：

\[
x_{t-\Delta t}
=
x_t
-
\Delta t\,v_\theta(x_t,t,x)
\]

把右侧记为一步均值：

\[
\mu_\theta(x_t,t,x)
=
x_t
-
\Delta t\,v_\theta(x_t,t,x)
\]

于是：

\[
x_{t-\Delta t}
=
\mu_\theta(x_t,t,x)
\]

这个式子说明：给定 $x_t$、时间 $t$ 和条件 $x$，下一步只有一个确定点。

如果硬写成概率分布：

\[
p_\theta(x_{t-\Delta t}\mid x_t,x,t)
=
\delta
\left(
x_{t-\Delta t}
-
\mu_\theta(x_t,t,x)
\right)
\]

这是 Dirac delta，不是普通高斯密度。

\begin{intubox}
\textbf{直觉：} ODE 像导航系统直接告诉你“下一步就走到这个位置”；RL 需要的是“下一步可以在这个位置附近随机采样，并知道每个采样点概率是多少”。
\end{intubox}

\subsubsection{第二步：ODE-to-SDE，给一步更新加随机性}

Flow-GRPO 的关键操作是把确定性 ODE transition 变成随机 SDE transition。

离散形式可以写成：

\[
x_{t-\Delta t}
=
\mu_\theta(x_t,t,x)
+
\sigma_t\sqrt{\Delta t}\,\epsilon
\]

其中：

\[
\epsilon\sim\mathcal{N}(0,I)
\]

也就是：

\[
x_{t-\Delta t}
\sim
\mathcal{N}
\left(
\mu_\theta(x_t,t,x),
\sigma_t^2\Delta t I
\right)
\]

于是每一步 flow transition 就有了普通高斯 logprob：

\[
\log p_\theta(x_{t-\Delta t}\mid x_t,x,t)
=
-
\frac{
\left\|
x_{t-\Delta t}
-
\mu_\theta(x_t,t,x)
\right\|^2
}{
2\sigma_t^2\Delta t
}
-
\frac{d}{2}
\log(2\pi\sigma_t^2\Delta t)
\]

这里 $d$ 是 latent 维度。

\begin{summarybox}
\textbf{关键变化：}\\
ODE：
\[
x_{t-\Delta t}=\mu_\theta
\]
没有好用 logprob。\\
SDE：
\[
x_{t-\Delta t}=\mu_\theta+\sigma_t\sqrt{\Delta t}\epsilon
\]
变成 Gaussian policy，因此可以计算 logprob。
\end{summarybox}

\subsubsection{第三步：整条 Flow Trajectory 的 Logprob}

一条 flow trajectory 写成：

\[
\tau
=
\left(
x_{t_K},x_{t_{K-1}},\ldots,x_{t_0}
\right)
\]

其中：

\[
t_K=1,
\qquad
t_0=0
\]

在 SDE 化后，每一步都有：

\[
p_\theta(x_{t_{k-1}}\mid x_{t_k},x,t_k)
\]

所以整条轨迹的 logprob 是：

\[
\log p_\theta(\tau\mid x)
=
\sum_{k=1}^{K}
\log
p_\theta(x_{t_{k-1}}\mid x_{t_k},x,t_k)
\]

这和 DDPO 完全对应：

\[
\log p_\theta(\tau_\text{diff}\mid x)
=
\sum_t
\log p_\theta(z_{t-1}\mid z_t,x,t)
\]

\begin{intubox}
\textbf{记忆法：} DDPO 是 diffusion step logprob 求和；Flow-GRPO 是 SDE flow step logprob 求和。
\end{intubox}

\subsubsection{第四步：Group Sampling}

GRPO 的核心不是训练 value model，而是对同一个条件生成一组候选，然后在组内做相对比较。

图像编辑条件写成：

\[
x=(I_\text{src},c)
\]

对同一个 $x$，采样 $G$ 条 flow trajectories：

\[
\tau_1,\tau_2,\ldots,\tau_G
\]

每条轨迹得到一个最终编辑图：

\[
I_j
=
D_\text{VAE}(x_{t_0}^{(j)})
\]

然后用 reward model / scorer 打分：

\[
R_j
=
R(I_\text{src},c,I_j)
\]

组内 advantage 是：

\[
A_j
=
\frac{
R_j-\text{mean}(R_1,\ldots,R_G)
}{
\text{std}(R_1,\ldots,R_G)+\epsilon
}
\]

\begin{summarybox}
\textbf{GRPO 的直觉：}\\
不需要知道某张图的绝对 reward 是否可靠，只要在同一个原图和同一个编辑指令下，知道哪几张图相对更好、哪几张图相对更差。
\end{summarybox}

\subsubsection{第五步：Flow-GRPO 的 Ratio}

采样时使用 old policy：

\[
\tau_j\sim p_{\theta_\text{old}}(\tau\mid x)
\]

对轨迹中的每一步，记录旧 logprob：

\[
\log p_{\theta_\text{old}}
\left(
x_{t_{k-1}}^{(j)}
\mid
x_{t_k}^{(j)},x,t_k
\right)
\]

更新时，用当前模型重新计算新 logprob：

\[
\log p_\theta
\left(
x_{t_{k-1}}^{(j)}
\mid
x_{t_k}^{(j)},x,t_k
\right)
\]

每一步 ratio 是：

\[
r_{j,k}(\theta)
=
\exp
\left(
\log p_\theta
\left(
x_{t_{k-1}}^{(j)}
\mid
x_{t_k}^{(j)},x,t_k
\right)
-
\log p_{\theta_\text{old}}
\left(
x_{t_{k-1}}^{(j)}
\mid
x_{t_k}^{(j)},x,t_k
\right)
\right)
\]

它的含义是：

\[
\text{同一个 flow transition，在新模型下比旧模型下更可能还是更不可能}
\]

Flow-GRPO 的 clipped objective 可以抽象写成：

\[
\mathcal{L}_\text{Flow-GRPO}
=
-
\mathbb{E}
\left[
\frac{1}{G}
\sum_{j=1}^{G}
\sum_{k=1}^{K}
\min
\left(
r_{j,k}(\theta)A_j,
\text{clip}(r_{j,k}(\theta),1-\epsilon,1+\epsilon)A_j
\right)
\right]
\]

如果 $A_j>0$，说明第 $j$ 个样本在组内比较好，训练会提高它的 trajectory probability。\\
如果 $A_j<0$，说明第 $j$ 个样本在组内比较差，训练会降低它的 trajectory probability。

\begin{intubox}
\textbf{直觉：} Flow-GRPO 不是直接告诉模型“应该怎么画”，而是告诉模型：“同一批候选里，高分轨迹更值得重复，低分轨迹更不值得重复”。
\end{intubox}

\subsubsection{图像编辑版 Flow-GRPO}

在 text-to-image 里，条件通常是：

\[
x=c_\text{text}
\]

在图像编辑里，条件变成：

\[
x=(I_\text{src},c)
\]

每一步 SDE transition 写成：

\[
p_\theta
\left(
x_{t_{k-1}}
\mid
x_{t_k},
I_\text{src},
c,
t_k
\right)
\]

最终 reward 必须同时评价：

\[
\text{编辑是否正确}
\quad
\text{原图是否保持}
\quad
\text{图像质量是否高}
\quad
\text{人类是否偏好}
\]

所以可以写成：

\[
R_j
=
w_\text{inst}R_\text{inst}^{(j)}
+
w_\text{pres}R_\text{pres}^{(j)}
+
w_\text{qual}R_\text{qual}^{(j)}
+
w_\text{pref}R_\text{pref}^{(j)}
\]

其中 preservation reward 尤其重要：

\[
R_\text{pres}^{(j)}
=
\cos
\left(
\phi(I_j),
\phi(I_\text{src})
\right)
\]

或者在 mask 外区域计算：

\[
R_\text{pres}^{(j)}
=
-
\text{LPIPS}
\left(
(1-m)\odot I_j,
(1-m)\odot I_\text{src}
\right)
\]

\begin{warnbox}
\textbf{编辑侧风险：} 如果 Flow-GRPO 只优化 PickScore / HPS / aesthetic reward，很容易把任务变成“重新生成一张漂亮图”。图像编辑必须有 source-image preservation 约束。
\end{warnbox}

\subsubsection{Flow-GRPO 为什么要 Denoising Reduction}

Flow-GRPO 的代价很高，因为每条 trajectory 有很多步：

\[
K\text{ steps}
\]

每个 condition 又要采样一组：

\[
G\text{ samples}
\]

如果每一步都算 logprob 和 ratio，训练成本大致随：

\[
G\times K
\]

增长。

所以 Flow-GRPO 提出 Denoising Reduction：训练时减少参与 RL 更新的 denoising / flow steps，但推理时仍保留原来的完整采样步数。

抽象地说，不再对所有时间步求和：

\[
\sum_{k=1}^{K}
\log p_\theta(x_{t_{k-1}}\mid x_{t_k},x,t_k)
\]

而是只选一个子集 $\mathcal{S}$：

\[
\sum_{k\in\mathcal{S}}
\log p_\theta(x_{t_{k-1}}\mid x_{t_k},x,t_k)
\]

这样能显著减少训练成本。

官方实现后续还加入了 Flow-GRPO-Fast：可以在 deterministic ODE trajectory 的某一个或两个中间 step 注入噪声，生成 group，只训练这些随机分叉点附近的 step。

\begin{intubox}
\textbf{直觉：} 完整 Flow-GRPO 是“整条轨迹都做 RL”；Denoising Reduction / Flow-GRPO-Fast 是“只挑关键几个步做 RL”，从而降低训练成本。
\end{intubox}

\subsubsection{Flow-GRPO 和 DDPO 的本质区别}

\begin{center}
{\small
\begin{tabular}{p{3cm}p{4.8cm}p{4.8cm}}
\hline
\textbf{问题} & \textbf{DDPO} & \textbf{Flow-GRPO} \\
\hline
底座模型
&
Diffusion model
&
Flow matching model
\\
原始采样
&
通常已经是随机反向过程
&
通常是 deterministic ODE
\\
logprob 来源
&
reverse Gaussian transition
&
ODE-to-SDE 后的 Gaussian transition
\\
advantage
&
reward baseline / PPO advantage
&
group-relative advantage
\\
是否需要 value model
&
可以有，也可以不用
&
通常不训练 value model
\\
主要工程难点
&
采样成本、reward 方差
&
SDE 化、step logprob、训练步数成本
\\
\hline
\end{tabular}
}
\end{center}

\begin{summarybox}
\textbf{最小区分：}\\
DDPO 的难点是“怎么把最终 reward 分给 diffusion steps”；Flow-GRPO 的额外难点是“Flow ODE 先天没有 step probability，必须先变成 SDE 才能做 ratio”。
\end{summarybox}

\subsubsection{LeapAlign：另一条 Flow 后训练路线}

Flow-GRPO 是 policy-gradient 路线：

\[
\text{sample}
\Rightarrow
\text{reward}
\Rightarrow
\text{logprob ratio update}
\]

LeapAlign 是 direct-gradient 路线：

\[
\text{reward}
\Rightarrow
\text{直接反传穿过 flow trajectory}
\]

问题是，完整 flow trajectory 很长。直接反传会带来：

\begin{itemize}[leftmargin=1.5em]
  \item 显存成本高；
  \item 梯度容易爆炸；
  \item 很难稳定更新早期生成步骤；
  \item 早期步骤往往决定整体布局和全局结构。
\end{itemize}

LeapAlign 的想法是：不要穿过完整长轨迹，而是构造 two-step leap trajectory。

假设从完整 ODE 轨迹中选两个时间点：

\[
k>j>0
\]

第一跳从 $x_k$ 直接预测到 $x_j$：

\[
\hat{x}_{j\mid k}
=
x_k
-
(k-j)
v_\theta(x_k,k,x)
\]

然后用 latent connector 把预测点接回真实轨迹点：

\[
x_j
=
\hat{x}_{j\mid k}
+
\text{stop\_gradient}
\left(
x_j-\hat{x}_{j\mid k}
\right)
\]

第二跳从 $x_j$ 直接预测到 clean 端：

\[
\hat{x}_{0\mid j}
=
x_j
-
j
v_\theta(x_j,j,x)
\]

再接回真实 $x_0$：

\[
x_0
=
\hat{x}_{0\mid j}
+
\text{stop\_gradient}
\left(
x_0-\hat{x}_{0\mid j}
\right)
\]

所以完整长路径被压缩成：

\[
x_k
\Rightarrow
\hat{x}_{j\mid k}
\Rightarrow
x_j
\Rightarrow
\hat{x}_{0\mid j}
\Rightarrow
x_0
\]

\begin{intubox}
\textbf{直觉：} LeapAlign 让模型学会“跳步预测”：不是一步一步走完整 ODE，而是从中间时刻直接跳到更靠近最终图像的位置。这样 reward gradient 只需要穿过两次 leap，而不是几十步采样。
\end{intubox}

\subsubsection{LeapAlign 的 Reward Objective}

LeapAlign 使用最终图像 reward：

\[
r(x_0)
\]

可以写一个 hinge-style loss：

\[
\mathcal{L}_\text{raw}
=
\max(0,\lambda-r(x_0))
\]

如果 reward 已经高于阈值 $\lambda$，loss 为 0；如果 reward 低于阈值，就推动模型提高 reward。

但不是所有 leap trajectory 都可靠。如果 leap 预测偏离真实长路径太远，反传信号就不可信。于是可以用 trajectory similarity 做权重。

例如：

\[
d_j
=
\text{mean}
\left(
|x_j-\hat{x}_{j\mid k}|
\right)
\]

\[
d_0
=
\text{mean}
\left(
|x_0-\hat{x}_{0\mid j}|
\right)
\]

定义：

\[
w_\text{sim}
=
\frac{1}
{\max(d_j,\tau)+\max(d_0,\tau)}
\]

最终 loss：

\[
\mathcal{L}
=
\text{stop\_gradient}(w_\text{sim})
\mathcal{L}_\text{raw}
\]

\begin{summarybox}
\textbf{LeapAlign 的核心：}\\
用 two-step leap trajectory 缩短反传路径，让 reward gradient 能更稳定地更新任意生成步骤，尤其是早期控制布局的步骤。
\end{summarybox}

\subsubsection{Flow-GRPO 和 LeapAlign 的关系}

\begin{center}
{\small
\begin{tabular}{p{3cm}p{4.8cm}p{4.8cm}}
\hline
\textbf{维度} & \textbf{Flow-GRPO} & \textbf{LeapAlign} \\
\hline
路线
&
Policy gradient / GRPO
&
Direct reward backpropagation
\\
是否需要 logprob
&
需要，所以要 ODE-to-SDE
&
不需要 logprob
\\
是否要求 reward 可微
&
不要求，只要能打分
&
更依赖可微 reward 或可反传 surrogate
\\
训练信号
&
group-relative advantage
&
reward gradient
\\
主要成本
&
采样多、ratio 计算多
&
反传轨迹显存和梯度稳定性
\\
核心技巧
&
SDE 化 + GRPO ratio + denoising reduction
&
two-step leap trajectory + similarity weighting + gradient discounting
\\
\hline
\end{tabular}
}
\end{center}

\begin{warnbox}
\textbf{不要混淆：} Flow-GRPO 和 LeapAlign 都是 flow matching 后训练，但不是同一个范式。Flow-GRPO 是“采样后用 reward 做 RL”；LeapAlign 是“让 reward 梯度直接穿过短轨迹反传”。
\end{warnbox}

\subsubsection{图像编辑里该怎么选}

如果 reward 是黑盒或离散 judge，例如：

\[
\text{VLM judge score}
\]

\[
\text{human preference scorer}
\]

\[
\text{OCR / object counting / rule-based reward}
\]

那么 Flow-GRPO 更自然，因为 policy gradient 不要求 reward 可微。

如果 reward 是可微的，例如某些 CLIP / HPS / differentiable aesthetic scorer，并且希望更高效地把梯度传给早期生成步骤，那么 LeapAlign 这类 direct-gradient 方法很有吸引力。

图像编辑任务里，实践上通常还要看：

\begin{itemize}[leftmargin=1.5em]
  \item 是否有可靠的 preservation reward；
  \item reward 是否可微；
  \item 训练预算是否允许 group sampling；
  \item 是否需要在线探索；
  \item 是否容易 reward hacking；
  \item 底座模型是 diffusion 还是 flow matching。
\end{itemize}

\begin{intubox}
\textbf{直觉：} Flow-GRPO 更像“用采样结果的好坏来调策略”；LeapAlign 更像“直接把评价器的梯度塞回生成过程”。前者适合黑盒 reward，后者适合可微 reward。
\end{intubox}

\subsubsection{本章最小结论}

\begin{summarybox}
\textbf{Flow-GRPO 的核心：}\\
Flow Matching 默认是 deterministic ODE，没有好用的 logprob。Flow-GRPO 先通过 ODE-to-SDE 把每一步 transition 变成 Gaussian policy，再对同一个条件采样 group，用 group-relative reward 构造 advantage，最后用新旧 step logprob ratio 做 GRPO 更新。\\
\textbf{LeapAlign 的核心：}\\
它不是 RL ratio 路线，而是 direct-gradient 路线。通过 two-step leap trajectory 缩短反传路径，让 reward gradient 能稳定更新任意 flow generation step，尤其是早期布局步骤。
\end{summarybox}

\subsubsection{本轮检查题}

\begin{enumerate}[leftmargin=1.5em]
  \item Flow Matching 的 ODE 为什么默认不能直接做 PPO / GRPO？
  \item ODE-to-SDE 后，每一步 transition 的 logprob 是怎么来的？
  \item Flow-GRPO 的 group-relative advantage 是什么？
  \item Flow-GRPO 的 ratio 和 DDPO 的 ratio 有什么对应关系？
  \item 图像编辑版 Flow-GRPO 为什么必须有 preservation reward？
  \item Denoising Reduction / Flow-GRPO-Fast 解决的是什么成本问题？
  \item LeapAlign 为什么不用 logprob？
  \item Flow-GRPO 和 LeapAlign 分别适合什么类型的 reward？
\end{enumerate}
